\section{Statistical inference and sensitivity}\label{sup:inference}

% The paper cites these subsections by number (S5.1 to S5.5): keep their order.

The paper reports one paragraph and one table of paired comparisons (its Section~5.10). This
section gives the rules behind them, every pair, and how the figures move when the rules
change. The scenarios were chosen by hand, several test the same mechanism (five are the
friction torque of a spinning conforming contact, three a body rolling against the foundation's
rolling resistance) and class B has five, so the bootstrap treats as exchangeable units that are
not: the intervals and the $p$-values are descriptive.

\subsection{Rules}

All summaries use the base variant. For two methods on the same scenarios we report the
geometric mean of the ratio of their errors with a 95\,\% percentile interval from 10\,000
bootstrap resamples within class, under the following rules.
\begin{itemize}
\item \emph{Cap and failures.} An error above \result{rules/cap} counts as \result{rules/cap},
  and so does a run without an error value, for whichever method it belongs to. No scenario is
  left out. Errors below \result{rules/floor} count as \result{rules/floor}. Without the cap, a
  run that ends with an enormous error (the two laws in \scn{massratio}, where the bodies fly
  apart) weighs several decades, by an amount that is arbitrary, while a run that fails without
  a value weighs nothing.
\item \emph{Clusters.} The four \scn{pieces} scenarios are the same problem on four cuts of
  the ground; they enter once, with the mean of their logarithmic ratios.
\item \emph{Ties.} A scenario is won or lost only if the two errors differ by more than a
  factor of \result{rules/tie}; otherwise it is a tie.
\item \emph{Leave-one-out.} Every ratio is given with the range that it covers when one
  unit is removed.
\item \emph{Primary comparisons and multiplicity.} \result{primary/n} comparisons are tested:
  \law{distributed} against \law{gonthier}, and each of the two against the default native
  model of each engine in this paper (Drake hydroelastic, Chrono SMC, MuJoCo soft), on class A
  and on classes A and B. Their bootstrap $p$-values, computed from the bootstrap distribution
  itself and not from a sample of it, are adjusted with Holm's method over that family, and
  ``significant'' means an adjusted $p$ below 0.05. Every other pair is exploratory and is
  reported without a test.
\end{itemize}
The cap, the tie band and the family were fixed after the data had been seen: the family is a
reduction of the number of tests, not a pre-registration. A law is compared with a native model
on that model's engine.

\subsection{Primary comparisons}

Table~\ref{tab:ratios} gives the primary comparisons with their leave-one-out ranges and their
counts of units won, tied and lost, next to the uncapped estimate, and Fig.~\ref{fig:forest}
their intervals. Of the \result{primary/n}, \result{primary/significant} are significant after
Holm's correction, both on classes A and B together: \law{distributed} against \law{gonthier}
(\result{robust/distributed-vs-gonthier/A+B/full}, adjusted $p = $
\result{robust/distributed-vs-gonthier/A+B/holm}) and against Chrono SMC
(\result{robust/distributed-vs-smc/A+B/full}, adjusted $p = $
\result{robust/distributed-vs-smc/A+B/holm}, close to the threshold). Two more have an
unadjusted interval that excludes 1 and do not survive the correction: \law{distributed}
against Drake hydroelastic (\result{robust/distributed-vs-hydro/A+B/full}) and against MuJoCo
soft (\result{robust/distributed-vs-soft/A+B/full}). On class A alone
\result{primary/significant/A} of the \result{primary/n/A} comparisons is significant: the two
laws are not separated from each other (\result{robust/distributed-vs-gonthier/A/full}, with
\result{robust/distributed-vs-gonthier/A/ties} ties in
\result{robust/distributed-vs-gonthier/A/n}), and against the default native models their
ratios are below 1 with intervals that reach above it (\law{distributed}:
\result{robust/distributed-vs-hydro/A/full} against Drake hydroelastic,
\result{robust/distributed-vs-smc/A/full} against Chrono SMC,
\result{robust/distributed-vs-soft/A/full} against MuJoCo soft). The count of scenarios says
more than the mean there: \law{distributed} has the smaller error by more than the tie band in
\result{robust/distributed-vs-hydro/A/wins}, \result{robust/distributed-vs-smc/A/wins} and
\result{robust/distributed-vs-soft/A/wins} of the \result{robust/distributed-vs-hydro/A/n}
units (won, tied and lost: \result{robust/distributed-vs-hydro/A/wtl},
\result{robust/distributed-vs-smc/A/wtl} and \result{robust/distributed-vs-soft/A/wtl}), and the
scenarios that it loses, it loses by much (\scn{massratio} against the two implicit solvers,
the sliding and rolling scenarios where a penalty or constraint model is at the floor).

\begin{table}[!htbp]
\caption{Primary comparisons: ratio of the errors of two methods on the same scenarios, base
variant. Reported: errors capped at \result{rules/cap}, runs without a value counted at the
cap, the \scn{pieces} scenarios one unit; $p$: bootstrap, Holm-adjusted over the
\result{primary/n} comparisons (bold: below 0.05); leave-one-out: range of the ratio when one
unit is removed; W/T/L: units won, tied (ratio within \result{rules/tie}) and lost by the first
method. Uncapped: raw errors, scenarios that a method did not finish left out (their number in
brackets), every scenario one unit}
\label{tab:ratios}
\generatedtable{ratios_primary}
\end{table}

\begin{figure}[!htbp]
\centering
\generatedfigure{forest_primary}
\caption{The primary comparisons with their 95\,\% bootstrap intervals, as reported and
uncapped}
\label{fig:forest}
\end{figure}

\subsection{Exploratory comparisons}

Table~\ref{tab:ratios-robust} gives every pair under the same rules, with unadjusted intervals
and no test. On class A \law{distributed} is not separated from Drake point
(\result{robust/distributed-vs-dpoint/A/full}) or from Chrono NSC
(\result{robust/distributed-vs-nsc/A/full}); on classes A and B its ratios to Drake point, to
Chrono NSC and to MuJoCo default are \result{robust/distributed-vs-dpoint/A+B/full},
\result{robust/distributed-vs-nsc/A+B/full} and
\result{robust/distributed-vs-mjdefault/A+B/full}. The error of \law{gonthier-thesis} is
\result{robust/thesis-vs-gonthier/A/full} times that of \law{gonthier} on class A: the friction
that Gonthier published, with the bristle on the material slip, is several times more accurate
than the variant of the 2018 thesis, which loses the scenarios in which a body must be held.
The error of \law{distributed} is \result{robust/distributed-vs-pcm/A+B/full} times that of
\law{pcm}.

\begin{table}[!htbp]
\caption{Every pair under the rules of the primary comparisons: geometric mean of the ratio of
the capped errors, 95\,\% bootstrap interval (unadjusted), and units won, tied and lost by the
first method}
\label{tab:ratios-robust}
\generatedtable{ratios_robust}
\end{table}

\subsection{Sensitivity analyses}

\emph{Without the cap.} Table~\ref{tab:ratios-first} gives the same pairs on the raw errors,
with the scenarios that either method did not finish left out and every scenario one unit; no
conclusion rests on it. On that estimate no interval of a law against a native model on class A
excludes 1, and the intervals span one to two orders of magnitude: \scn{massratio} enters with
the nearly seven decades between the error of a lost run and the floor, while the runs that the
native models did not finish are left out. Capping the errors and counting the failures moves
every primary ratio of a law against a native model towards the laws (on class A, from between
\result{ratio/distributed-vs-hydro/A/gm} and \result{ratio/gonthier-vs-soft/A/gm} to between
\result{robust/distributed-vs-hydro/A/gm} and \result{robust/gonthier-vs-soft/A/gm}) and
narrows the intervals; the choice of rules therefore favours the laws. For \law{gonthier-thesis} against
\law{gonthier} the uncapped ratio on class A is \result{ratio/thesis-vs-gonthier/A/full}.

\emph{Leave-one-out.} Removing any one unit leaves every primary ratio of a law against a
native model below or at 1 (Table~\ref{tab:ratios}).

\emph{Reduced class A.} Table~\ref{tab:reduced} repeats the primary comparisons on class A
without \scn{bowl}, \scn{trough}, \scn{turntable} and \scn{edge}, whose references no damped
compliant model can meet at the $h$ and $K$ of the battery: the ratios move by less than their
intervals, and none excludes 1. The counts within 5\,\% on the reduced class are
\result{pass/Ar/distributed/5} for \law{distributed}, \result{pass/Ar/gonthier/5} for
\law{gonthier}, and between \result{passn/Ar/mjdefault/5} and \result{passn/Ar/hydro/5} for the
native models (on the whole class A, \result{pass/A/distributed/5},
\result{pass/A/gonthier/5}, and between \result{passn/A/mjdefault/5} and
\result{passn/A/hydro/5}; within 1\,\%, \result{passn/A/distributed/1} for \law{distributed}
and \result{passn/A/dpoint/1} for Drake point; on class B within 5\,\%,
\result{pass/B/distributed/5} for \law{distributed}, \result{pass/B/hydro/5} for Drake
hydroelastic, \result{pass/B/gonthier/5} for \law{gonthier} and at most
\result{passn/B/smc/5} for a point-contact model).

\emph{Perturbed starts.} Over the eight perturbations of the initial velocities the range of
the error relative to its median has a median over the scenarios of \result{pert/gonthier} for
\law{gonthier} and \result{pert/distributed} for \law{distributed}: the base errors of the laws
are representative of nearby initial conditions. Those of several native models are not
(Section~\ref{sup:scatter}), so for them single base runs enter the ratios as draws.

\begin{table}[!htbp]
\caption{Uncapped estimate: ratio of the raw errors of two methods on the same scenarios, base
variant: geometric mean, 95\,\% bootstrap interval, and scenarios won and lost by the first
method (in brackets, the scenarios left out because one of the two did not finish). Bold: the
unadjusted interval excludes 1}
\label{tab:ratios-first}
\generatedtable{ratios}
\end{table}

\begin{table}[!htbp]
\caption{The primary comparisons on class A and on class A without the four parameter-bound
scenarios}
\label{tab:reduced}
\generatedtable{reduced}
\end{table}

\subsection{Runs that would settle the open points}

The following would answer open points of the paper and were not done. (1) Log, for every run,
whether the step bounds were active and the extreme ratios of the applied to the nominal
parameters ($h$, $d$, $v_s$, stick stiffness). (2) Record the sign of the class B errors, to
place every run against the band (the paper's Section~5.4 signs those of \scn{disk} and
\scn{bearing} only). (3) One more decade of stiffness ($K\times100$) on class A,
at the step that the laws need. (4) \scn{journal} at 256 segments with the climb angle
reported, and \scn{cone} at $10K$ with 128 segments, to separate the facet sag from the
self-excited wobble. (5) Sweeps of $d$ and $\tau$ on the scenarios that hold a body
(\scn{push}, \scn{wedge}, \scn{tower}, \scn{stack}, \scn{tip}, \scn{hole}), and of $h$ on the
rolling ones. (6) Native models again with reviewed configurations: Drake point with rim
spheres on cylinder faces, MuJoCo soft on \scn{block}, repeated or better converged runs of
Chrono NSC, and a check of the static penetration of every native model per scenario. (7) A
bisection of the change of \law{gonthier} in \scn{energy} between the first run of the second
set and the battery, and an energy measure that counts the contact spring. (8) The linear
stability model with the discrete sums of \law{pcm}, to explain why it does not flutter. (9)
The quadratures on remeshed flat faces inside the library, with their cost at equal accuracy.
(10) Finite-element solutions for the other class B scenarios (\scn{ring}, \scn{contensou},
\scn{cone}; the paper's Section~5.4 covers \scn{disk} and \scn{bearing}, statically), a
dynamic one with friction, measured data for one class B case, and three to five scenarios
written and frozen before any run, as a true out-of-sample check.
